% Section 4 insert generated from the executed deterministic replay.
% Production reference: 010518fbece77a4630684065101033ff8588613d
\subsection{Deterministic replay of machine-specific diagnostic policies}
\label{sec:policy_replay}

We evaluated machine-aware diagnostic policies by replaying eight fixed,
three-sample sequences through the production policy classes without
communicating with EPICS.  All inputs were synthetic.  They comprised two
nominal finite sequences, a constant-zero ambiguity control, a constant
nonzero proxy for a frozen readback, a sustained high-level signal, an isolated
$10^8$-level spike, and sequences containing NaN or infinity.  The spike
amplitude reflects the reported order of magnitude of anomalous IRFEL
energy-meter readings, but the sequence was constructed rather than recorded
on the machine.  For each sequence, we recorded the arithmetic mean, mean and
maximum absolute amplitudes, peak-to-peak variation, processed objective, and
events emitted by the production policy event sink.

Three objective configurations were compared: no policy; the dedicated
\texttt{FelEnergyGuardPolicy}; and a registry-built
\texttt{SampleGuardObjectivePolicy}.  The dedicated guard replaced the reduced
target with zero when the mean absolute sample amplitude exceeded $10^6$ or
when peak-to-peak variation was below $10^{-6}$.  The declarative rule used the
named target \texttt{fel\_energy}, the conditions
\texttt{max\_abs > $10^6$} or \texttt{peak\_to\_peak < $10^{-6}$}, and a
\texttt{replace: 0} action.  These configurations are not asserted to be
equivalent for arbitrary inputs; they are compared on the fixed replay.  A
separate constraint replay configured a named \texttt{orbit\_x} rule to map an
all-zero BPM sequence outside the accepted interval $[-1,1]$.

\begin{table}[t]
\centering
\caption{Observed results for eight synthetic objective sequences.  ``Changed''
means that the post-policy scalar differed from the arithmetic mean; it is not
a classification-accuracy measure.}
\label{tab:policy-replay}
\begin{tabular}{lrrrr}
\hline
Configuration & Changed & Finite & Nominal retained & Observable events \\
\hline
No policy & 0/8 & 6/8 & 2/2 & 0 \\
FEL energy guard & 4/8 & 7/8 & 2/2 & 0$^{\ast}$ \\
Declarative sample guard & 4/8 & 7/8 & 2/2 & 5 \\
\hline
\end{tabular}
\\[-1mm]
\footnotesize{$^{\ast}$The production FEL guard does not emit policy events; zero
events therefore denotes unavailable event observability, not a false
predicate.}
\end{table}

Both guards mapped the isolated $[2,10^8,2]$ sequence from its arithmetic mean
of $3.3333334667\times10^7$ to zero, and both changed all three constructed
anomaly sequences while preserving both nominal finite sequences.  Thus the
replay does not show an isolated-spike defect in the dedicated FEL guard and
does not establish superiority of the general rule.  Instead, it demonstrates
that the same optimizer-facing evaluation interface can host either a
machine-specific policy class or a named declarative rule.  The declarative
guard emitted five observable events, including one for the constant-zero
control whose numerical value remained zero.  The FEL guard emitted no events
because its production implementation has no event-emission call; no internal
trigger log was inferred from its outputs.

Without a policy, NaN and infinity both propagated.  Under either guard, NaN
still propagated, whereas infinity was replaced because it satisfied the
amplitude predicate.  This behavior should not be interpreted as comprehensive
nonfinite-data handling.  In the constraint replay, the normal BPM mean
remained $0.3$, while the all-zero sequence was mapped to $1.2$, deliberately
outside $[-1,1]$, and produced one policy event.  The experiment exercised
\texttt{post\_reduce()} only; it did not test PV acquisition or writes,
snapshot/restore, GUI transitions, optimizer convergence, or live-machine
safety.  Frozen-value discrimination also remains intrinsically ambiguous
when a physically valid signal can be constant.
